% GENERATED FILE. Edit canonical lessons, then run npm run generate:books.
% canonical-source-hash: fnv1a64:10c0e3b3cae21b6d
% canonical-lessons: SA-C463-bhisayate, SA-C463-vihvalah, SA-C463-bhiruh, SA-C463-trasyati, SA-C463-pragalbhah, SA-R463-review-456-463

\chapter{Courage and Fear}
\label{ch:sa-courage-and-fear}
\hlchaptermodality{463}

\begin{chapteropening}
\textbf{By the end of this chapter:} \emph{I can say whether someone is brave or timid.}

Complete the last lesson of chapter 463: I can say whether someone is brave or timid.
\end{chapteropening}

\section[\texorpdfstring{\sk{भीषयते} (bhīṣayate)}{भीषयते (bhīṣayate)}]{\sk{भीषयते} (bhīṣayate) — frightens}

\label{lesson:SA-C463-bhisayate}



\begin{quote}
Before the new one: say the Sanskrit for trustworthy, then the Sanskrit for a betrayal of trust.
\end{quote}

\subsection*{You'll want to know: \sk{भीषयते}}
\textbf{\sk{भीषयते}} — \emph{bhīṣayate} — \textquotedblleft{}frightens\textquotedblright{}.

\begin{quote}
\sk{सिंहः} \sk{गर्जनेन} \sk{मृगान्} \sk{भीषयते।}
\end{quote}

\begin{quote}
\emph{siṃhaḥ garjanena mṛgān bhīṣayate.}
\end{quote}

\begin{quote}
The lion frightens the deer with its roar.
\end{quote}

Past (it …): \sk{सिंहः} \sk{मृगान्} \sk{भीषितवान्}, \emph{siṃhaḥ mṛgān bhīṣitavān}.

\begin{cousinweb}
\textbf{\sk{भीषयते}} (\emph{bhīṣayate}) is built on the causative of the root √bhī, \textquotedblleft{}to fear\textquotedblright{}, the root of \textbf{\sk{बिभेति}} and \textbf{\sk{भयम्}}: to make afraid; MW records this causative \emph{bhīṣayate}, citing Pāṇini 1.3.68 (MW).
\end{cousinweb}

\begin{grammarlens}[title={What you've built}]
One more word for describing what you lived and felt.
\end{grammarlens}

\subsection*{Your turn}
\begin{itemize}
\raggedright
  \item \emph{Say it:} \emph{bhīṣayate}
  \item \emph{Say it:} \emph{bhīṣayate}, once more
  \item \emph{Say it:} \emph{siṃhaḥ garjanena mṛgān bhīṣayate.} — The lion frightens the deer with its roar.
  \item \emph{Recall:} say the Sanskrit for trustworthy, then the Sanskrit for a betrayal of trust, then say \emph{bhīṣayate} again
\end{itemize}

\begin{tcolorbox}[breakable,colback=teal!4,colframe=teal!35!black,title={Before you move on}]
What does \sk{भीषयते} mean? (\textquotedblleft{}Frightens\textquotedblright{}.) Say it once more.
\end{tcolorbox}
\section[\texorpdfstring{\sk{विह्वलः} (vihvalaḥ)}{विह्वलः (vihvalaḥ)}]{\sk{विह्वलः} (vihvalaḥ) — distraught; perturbed}

\label{lesson:SA-C463-vihvalah}



\begin{quote}
Before the new one: say the Sanskrit for a betrayal of trust, then the Sanskrit for frightens.
\end{quote}

\subsection*{You'll want to know: \sk{विह्वलः}}
\textbf{\sk{विह्वलः}} — \emph{vihvalaḥ} — \textquotedblleft{}distraught; perturbed\textquotedblright{}.

\begin{quote}
\sk{भयेन} \sk{विह्वलः} \sk{बालः} \sk{मातरम्} \sk{आह्वयति।}
\end{quote}

\begin{quote}
\emph{bhayena vihvalaḥ bālaḥ mātaram āhvayati.}
\end{quote}

\begin{quote}
Distraught with fear, the child calls his mother.
\end{quote}

\begin{cousinweb}
\textbf{\sk{विह्वलः}} (\emph{vihvalaḥ}) is built from \emph{vi} and the root √hval, \textquotedblleft{}to go crookedly, to stagger\textquotedblright{}; MW gives \textquotedblleft{}agitated, perturbed, distressed, afflicted\textquotedblright{} (MW).
\end{cousinweb}

\begin{grammarlens}[title={What you've built}]
One more word for describing what you lived and felt.
\end{grammarlens}

\subsection*{Your turn}
\begin{itemize}
\raggedright
  \item \emph{Say it:} \emph{vihvalaḥ}
  \item \emph{Say it:} \emph{vihvalaḥ}, once more
  \item \emph{Say it:} \emph{bhayena vihvalaḥ bālaḥ mātaram āhvayati.} — Distraught with fear, the child calls his mother.
  \item \emph{Recall:} say the Sanskrit for a betrayal of trust, then the Sanskrit for frightens, then say \emph{vihvalaḥ} again
\end{itemize}

\begin{tcolorbox}[breakable,colback=teal!4,colframe=teal!35!black,title={Before you move on}]
What does \sk{विह्वलः} mean? (\textquotedblleft{}Distraught; perturbed\textquotedblright{}.) Say it once more.
\end{tcolorbox}
\section[\texorpdfstring{\sk{भीरुः} (bhīruḥ)}{भीरुः (bhīruḥ)}]{\sk{भीरुः} (bhīruḥ) — timid}

\label{lesson:SA-C463-bhiruh}



\begin{quote}
Before the new one: say the Sanskrit for frightens, then the Sanskrit for distraught.
\end{quote}

\subsection*{You'll want to know: \sk{भीरुः}}
\textbf{\sk{भीरुः}} — \emph{bhīruḥ} — \textquotedblleft{}timid\textquotedblright{}.

\begin{quote}
\sk{भीरुः} \sk{शशकः} \sk{सहसा} \sk{धावति।}
\end{quote}

\begin{quote}
\emph{bhīruḥ śaśakaḥ sahasā dhāvati.}
\end{quote}

\begin{quote}
The timid rabbit suddenly runs off.
\end{quote}

\begin{cousinweb}
\textbf{\sk{भीरुः}} (\emph{bhīruḥ}) comes from the root √bhī, \textquotedblleft{}to fear\textquotedblright{}, the root of \textbf{\sk{बिभेति}} and \textbf{\sk{भयम्}}; MW gives \textquotedblleft{}fearful, timid, cowardly\textquotedblright{} (MW).
\end{cousinweb}

\begin{grammarlens}[title={What you've built}]
One more word for describing what you lived and felt.
\end{grammarlens}

\subsection*{Your turn}
\begin{itemize}
\raggedright
  \item \emph{Say it:} \emph{bhīruḥ}
  \item \emph{Say it:} \emph{bhīruḥ}, once more
  \item \emph{Say it:} \emph{bhīruḥ śaśakaḥ sahasā dhāvati.} — The timid rabbit suddenly runs off.
  \item \emph{Recall:} say the Sanskrit for frightens, then the Sanskrit for distraught, then say \emph{bhīruḥ} again
\end{itemize}

\begin{tcolorbox}[breakable,colback=teal!4,colframe=teal!35!black,title={Before you move on}]
What does \sk{भीरुः} mean? (\textquotedblleft{}Timid\textquotedblright{}.) Say it once more.
\end{tcolorbox}
\section[\texorpdfstring{\sk{त्रस्यति} (trasyati)}{त्रस्यति (trasyati)}]{\sk{त्रस्यति} (trasyati) — trembles with fear}

\label{lesson:SA-C463-trasyati}



\begin{quote}
Before the new one: say the Sanskrit for distraught, then the Sanskrit for timid.
\end{quote}

\subsection*{You'll want to know: \sk{त्रस्यति}}
\textbf{\sk{त्रस्यति}} — \emph{trasyati} — \textquotedblleft{}trembles with fear\textquotedblright{}.

\begin{quote}
\sk{मूषकः} \sk{मार्जारात्} \sk{त्रस्यति।}
\end{quote}

\begin{quote}
\emph{mūṣakaḥ mārjārāt trasyati.}
\end{quote}

\begin{quote}
The mouse trembles in fear of the cat.
\end{quote}

Past (it …): \sk{मूषकः} \sk{मार्जारात्} \sk{त्रस्तवान्}, \emph{mūṣakaḥ mārjārāt trastavān}.

\begin{cousinweb}
\textbf{\sk{त्रस्यति}} (\emph{trasyati}) comes from the root √tras, \textquotedblleft{}to tremble, to be afraid of\textquotedblright{}, which takes the ablative of what is feared (MW).
\end{cousinweb}

\begin{grammarlens}[title={What you've built}]
One more word for describing what you lived and felt.
\end{grammarlens}

\subsection*{Your turn}
\begin{itemize}
\raggedright
  \item \emph{Say it:} \emph{trasyati}
  \item \emph{Say it:} \emph{trasyati}, once more
  \item \emph{Say it:} \emph{mūṣakaḥ mārjārāt trasyati.} — The mouse trembles in fear of the cat.
  \item \emph{Recall:} say the Sanskrit for distraught, then the Sanskrit for timid, then say \emph{trasyati} again
\end{itemize}

\begin{tcolorbox}[breakable,colback=teal!4,colframe=teal!35!black,title={Before you move on}]
What does \sk{त्रस्यति} mean? (\textquotedblleft{}Trembles with fear\textquotedblright{}.) Say it once more.
\end{tcolorbox}
\section[\texorpdfstring{\sk{प्रगल्भः} (pragalbhaḥ)}{प्रगल्भः (pragalbhaḥ)}]{\sk{प्रगल्भः} (pragalbhaḥ) — bold}

\label{lesson:SA-C463-pragalbhah}



\begin{quote}
Before the new one: say the Sanskrit for timid, then the Sanskrit for trembles with fear.
\end{quote}

\subsection*{You'll want to know: \sk{प्रगल्भः}}
\textbf{\sk{प्रगल्भः}} — \emph{pragalbhaḥ} — \textquotedblleft{}bold\textquotedblright{}.

\begin{quote}
\sk{प्रगल्भः} \sk{बालः} \sk{सभायां} \sk{वदति।}
\end{quote}

\begin{quote}
\emph{pragalbhaḥ bālaḥ sabhāyāṃ vadati.}
\end{quote}

\begin{quote}
The bold boy speaks up in the assembly.
\end{quote}

\begin{cousinweb}
\textbf{\sk{प्रगल्भः}} (\emph{pragalbhaḥ}) is built from \emph{pra} and the root √galbh, \textquotedblleft{}to be bold\textquotedblright{}; MW gives \textquotedblleft{}bold, confident, resolute, brave\textquotedblright{} (MW).
\end{cousinweb}

\begin{grammarlens}[title={What you've built}]
One more word for describing what you lived and felt.
\end{grammarlens}

\subsection*{Your turn}
\begin{itemize}
\raggedright
  \item \emph{Say it:} \emph{pragalbhaḥ}
  \item \emph{Say it:} \emph{pragalbhaḥ}, once more
  \item \emph{Say it:} \emph{pragalbhaḥ bālaḥ sabhāyāṃ vadati.} — The bold boy speaks up in the assembly.
  \item \emph{Recall:} say the Sanskrit for timid, then the Sanskrit for trembles with fear, then say \emph{pragalbhaḥ} again
\end{itemize}

\begin{tcolorbox}[breakable,colback=teal!4,colframe=teal!35!black,title={Before you move on}]
What does \sk{प्रगल्भः} mean? (\textquotedblleft{}Bold\textquotedblright{}.) Say it once more.
\end{tcolorbox}
\section[(dialogue)]{Looking back: from Calm and Composure to Courage and Fear}

\label{lesson:SA-R463-review-456-463}



\begin{quote}
Say the Sanskrit for is shaken and appeases.
\end{quote}

\subsection*{Your turn}
Say these aloud:

\begin{itemize}
\raggedright
  \item is shaken, appeases, composed, massages, a secluded place
  \item smiles, laughs at, a jest, a loud laugh, amusement
  \item haughtiness, boasts, bombast, arrogant, modest
  \item contempt, despises, insults, regard, self-respect
  \item temperament, refined, conducts himself, a virtue, wicked
  \item a great soul, a scoundrel, goodness, heartless, bears malice
  \item deceitful, devious, explains away, trustworthy, a betrayal of trust
  \item frightens, distraught, timid, trembles with fear, bold
\end{itemize}

\begin{tcolorbox}[breakable,colback=teal!4,colframe=teal!35!black,title={Before you move on}]
Is shaken? (\textbf{\sk{क्षुभ्यति}}, \emph{kṣubhyati}.) Appeases? (\textbf{\sk{प्रसादयति}}, \emph{prasādayati}.) Composed? (\textbf{\sk{समाहितः}}, \emph{samāhitaḥ}.) Massages? (\textbf{\sk{मर्दयति}}, \emph{mardayati}.) A secluded place? (\textbf{\sk{एकान्तः}}, \emph{ekāntaḥ}.) Smiles? (\textbf{\sk{स्मयते}}, \emph{smayate}.) Laughs at? (\textbf{\sk{उपहसति}}, \emph{upahasati}.) A jest? (\textbf{\sk{परिहासः}}, \emph{parihāsaḥ}.) A loud laugh? (\textbf{\sk{अट्टहासः}}, \emph{aṭṭahāsaḥ}.) Amusement? (\textbf{\sk{विनोदः}}, \emph{vinodaḥ}.) Haughtiness? (\textbf{\sk{दर्पः}}, \emph{darpaḥ}.) Boasts? (\textbf{\sk{श्लाघते}}, \emph{ślāghate}.) Bombast? (\textbf{\sk{आडम्बरः}}, \emph{āḍambaraḥ}.) Arrogant? (\textbf{\sk{उद्धतः}}, \emph{uddhataḥ}.) Modest? (\textbf{\sk{विनीतः}}, \emph{vinītaḥ}.) Contempt? (\textbf{\sk{अवज्ञा}}, \emph{avajñā}.) Despises? (\textbf{\sk{अवमन्यते}}, \emph{avamanyate}.) Insults? (\textbf{\sk{तिरस्करोति}}, \emph{tiraskaroti}.) Regard? (\textbf{\sk{आदरः}}, \emph{ādaraḥ}.) Self-respect? (\textbf{\sk{मानः}}, \emph{mānaḥ}.) Temperament? (\textbf{\sk{स्वभावः}}, \emph{svabhāvaḥ}.) Refined? (\textbf{\sk{सभ्यः}}, \emph{sabhyaḥ}.) Conducts himself? (\textbf{\sk{प्रवर्तते}}, \emph{pravartate}.) A virtue? (\textbf{\sk{गुणः}}, \emph{guṇaḥ}.) Wicked? (\textbf{\sk{दुष्टः}}, \emph{duṣṭaḥ}.) A great soul? (\textbf{\sk{महात्मा}}, \emph{mahātmā}.) A scoundrel? (\textbf{\sk{खलः}}, \emph{khalaḥ}.) Goodness? (\textbf{\sk{सौजन्यम्}}, \emph{saujanyam}.) Heartless? (\textbf{\sk{निष्ठुरः}}, \emph{niṣṭhuraḥ}.) Bears malice? (\textbf{\sk{द्रुह्यति}}, \emph{druhyati}.) Deceitful? (\textbf{\sk{शठः}}, \emph{śaṭhaḥ}.) Devious? (\textbf{\sk{कुटिलः}}, \emph{kuṭilaḥ}.) Explains away? (\textbf{\sk{अपलपति}}, \emph{apalapati}.) Trustworthy? (\textbf{\sk{विश्वसनीयः}}, \emph{viśvasanīyaḥ}.) A betrayal of trust? (\textbf{\sk{विश्वासघातः}}, \emph{viśvāsaghātaḥ}.) Frightens? (\textbf{\sk{भीषयते}}, \emph{bhīṣayate}.) Distraught? (\textbf{\sk{विह्वलः}}, \emph{vihvalaḥ}.) Timid? (\textbf{\sk{भीरुः}}, \emph{bhīruḥ}.) Trembles with fear? (\textbf{\sk{त्रस्यति}}, \emph{trasyati}.) Bold? (\textbf{\sk{प्रगल्भः}}, \emph{pragalbhaḥ}.)
\end{tcolorbox}
